\section{引言}
\label{sec:intro}

截至 2026 年，大语言模型（LLMs）的演进已形成泾渭分明的两极。一方面，如 GPT-5.3、Claude 4.6 Opus 和 Gemini 3.1 Pro 这样的旗舰模型，在逻辑推理和复杂工具调用（Agentic Coding）上展现出了令人惊叹的能力；但另一方面，它们高昂的 API 调用成本使得多智能体（Multi-Agent）长时间、大规模的“左右互搏”变得极其不经济。

与之对应的是，开源社区和下沉市场充斥着大量价格极其低廉的模型（如 MiniMax M2.5 或各路量化架构）。这些模型虽然在文本生成上速度极快，但在作为独立智能体执行任务时，通常面临着灾难性的缺陷：
\begin{enumerate}
    \item \textbf{JSON 恐惧症与工具失控}：它们极难严格遵循复杂的结构化 Tool-Calling 规范，稍微复杂的 ReAct 循环或沙盒 API 就会让它们逻辑崩溃，这也导致直接赋予它们 OpenClaw 级别的物理环境权限变得极其危险。
    \item \textbf{金鱼般的记忆力}：在多轮辩论中，便宜模型的长上下文衰减极其严重，它们经常在聊到第 10 轮时就忘了第一轮的核心假设，开始车轱辘话来回说。
    \item \textbf{幻觉传染}：由于缺乏底层逻辑定力，一旦某个智能体在一本正经地胡说八道（如捏造不存在的文献），其他廉价模型不仅不会反驳，反而极易陷入“盲从”，形成可怕的“幻觉回音室（Echo Chamber）”。
\end{enumerate}

如何将这群“不靠谱”的廉价模型压榨出科研价值？本文提出了 \textbf{Grox\_chat} 框架。本框架的核心哲学可以概括为一句带有黑色幽默的宣言：“不要指望垃圾模型拥有自律，只能用最严酷的物理监狱去牵制它们。”

在这个名为 Grox\_chat 的监狱式科研系统中，我们利用了顶级模型（Gemini 3.1 Pro）担任高高在上的典狱长（负责宏观调度与事实核查），而将高强度的辩论脏活累活全部下放给廉价模型。为了让廉价模型乖乖就范，我们在架构上做出了三项“降维打击”式的贡献：
\begin{itemize}
    \item \textbf{剥夺记忆权：持久化黑板与总结查重}。模型不再被允许自己维持上下文，所有发言都被强制写入 SQLite 数据库。典狱长通过定期的 Summary RAG 来进行“查重”，一旦发现专家们在水时长或陷入循环，立刻强行拉闸终止话题。
    \item \textbf{剥夺工具权：自然语言意图 RAG}。彻底废弃容易引发系统崩溃的 JSON API 协议。廉价模型只需说一句“帮我查一下某某理论”，底层系统会自动截获这句大白话，并使用本地 \texttt{sqlite-vec} 与交叉编码器去事实库里提取硬核证据贴回它的脸上。
    \item \textbf{引入疯狗机制：DAPA 异步制裁}。在正常的专家辩论序列之外，我们在 LangGraph 中散养了三只名为“猫”、“狗”和“Tron”的异步守护程序。它们不讲武德，一旦发现有人造假或逻辑断裂，狗（Dot）就会在当轮回合末尾强行插入惩罚序列，逼迫造假者重写论点。
\end{itemize}

Grox\_chat 证明了，即使基座模型破绽百出，只要架构的枷锁足够沉重，依然能让它们在沙盒里推演出极高质量的科研闭环。